% =========================================================================
% MLC-Norm — final project paper, Machine Learning Compilers Lab, FSU Jena.
%
% Build:   pdflatex main && pdflatex main      (twice, for cross-references)
%   or:    latexmk -pdf main.tex
%
% Dependencies, all standard in TeX Live / Overleaf:
%   IEEEtran, booktabs, amsmath, graphicx, hyperref, url
%
% Figures are PDFs under paper/figures/, generated from the same sources as
% the HTML report's SVGs:
%   python3 docs/figures/make_figures.py   # data -> SVG (docs/source/_static)
%   python3 docs/figures/svg2pdf.py        # SVG  -> PDF (paper/figures)
%
% If IEEEtran is unavailable, replace the \documentclass line with
%   \documentclass[10pt,twocolumn]{article}
% and delete \IEEEpeerreviewmaketitle; everything else compiles unchanged.
% =========================================================================
\documentclass[conference,10pt]{IEEEtran}

\usepackage[T1]{fontenc}
\usepackage[utf8]{inputenc}
\usepackage{amsmath}
\usepackage{amssymb}
\usepackage{booktabs}
\usepackage{graphicx}
\usepackage{url}
\usepackage[hidelinks]{hyperref}

% Compact code style without extra package dependencies.
\newcommand{\code}[1]{\texttt{\small #1}}

\begin{document}

\title{MLC-Norm: Hand-Written and JIT-Generated SME/SSVE\\
       Kernels for LayerNorm and RMSNorm on Apple Silicon}

\author{%
\IEEEauthorblockN{Ketsia Kemkuini and Mariza Yamdjeu}
\IEEEauthorblockA{Machine Learning Compilers Lab\\
Friedrich-Schiller-Universit\"at Jena\\
\url{https://github.com/kets01/MLC-Project}}
}

\maketitle

\begin{abstract}
Transformer normalization runs in every block, twice per layer, and touches
every element of the activation tensor while doing very little arithmetic per
element. It is therefore a data-movement problem. We ask how close hand-written
and JIT-generated SME/Streaming-SVE kernels can bring LayerNorm and RMSNorm to
the achievable memory bandwidth of an Apple M4, and---more usefully---what
actually limits them.

We take both norms from a scalar reference through seven vectorized variants, a
Scalable Matrix Extension (ZA) variant, runtime code generation, and integration
into a tiled-execution-IR runtime, verifying every configuration against a
float64 oracle before it is timed. The best kernels reach 33\,\% of a
footprint-matched bandwidth ceiling on useful bytes at cache-resident shapes and
42\,\% in the DRAM regime, and outrun PyTorch 2.13 and ExecuTorch 1.4.1 by
1.5--2.1$\times$ (LayerNorm) and 4.2--5.8$\times$ (RMSNorm) single-threaded.

The negative and corrective results are the more transferable ones. Using the
ZA matrix accumulator as a residency buffer \emph{loses} by 24--37\,\% even
after an SME2 rebuild that made it 62--114\,\% faster, because it trades a
memory read that was never a memory read for \code{mova} traffic. The
streaming-mode transition costs 9.04\,ns, roughly a fifth of the per-call floor,
so the small-tensor penalty is kernel group granularity rather than
\code{SMSTART}. And our own roofline was mis-stated for five sprints: the
bandwidth ceiling is a curve in working-set size, not a constant, and dividing
cache-resident kernels by the DRAM figure had inflated our reported efficiency
by about 2$\times$.
\end{abstract}

\begin{IEEEkeywords}
SME, Streaming SVE, Apple silicon, LayerNorm, RMSNorm, roofline, JIT code
generation, memory bandwidth.
\end{IEEEkeywords}

\IEEEpeerreviewmaketitle

% =========================================================================
\section{Introduction}

Normalization is not an afterthought in a Transformer. A norm runs in every
block, twice per layer, of the Transformer architecture~\cite{vaswani2017attention},
in models built on LayerNorm~\cite{ba2016layernorm} (BERT, GPT-2/3) and on
RMSNorm~\cite{zhang2019rmsnorm} (LLaMA~\cite{touvron2023llama}, Gemma, Mistral).
Both read and write an entire activation tensor while performing only a handful
of floating-point operations per element, which places them firmly on the
bandwidth-dominated side of the roofline~\cite{williams2009roofline} at the
shapes we measure. Making them fast is a data-movement problem.

This paper reports an engineering study of that problem on one concrete target:
the Arm Scalable Matrix Extension (SME) and its Streaming SVE (SSVE) mode as
implemented on the Apple M4. The question we set out to answer has two halves:

\begin{quote}
\emph{How close to the machine's achievable memory bandwidth can hand-written
and JIT-generated SME/SSVE kernels bring LayerNorm and RMSNorm, and what
actually limits them?}
\end{quote}

The second half proved more valuable than the first. A GiB/s number is only
meaningful against a correctly chosen ceiling, and a substantial part of what
this project learned came from discovering that its own ceiling---and several of
its own published claims---were wrong.

\subsection*{Contributions}

\begin{enumerate}
\item \textbf{Two kernel families, fully ablated.} Seven SSVE variants and two
ZA variants per norm, each isolating one lever, each verified against a float64
reference before any throughput figure is produced (Sec.~\ref{sec:ablation}).
The winning lever is \emph{access density}, not the residency we expected.

\item \textbf{A measured negative result for the ZA tile.} Staging activations
through the SME matrix accumulator to eliminate a pass loses decisively, and
still loses after an SME2 rebuild that made the staging 4$\times$ cheaper. We
report the mechanism and the retraction of our own earlier reasoning
(Sec.~\ref{sec:za}).

\item \textbf{A correction to our own roofline.} The single-core bandwidth
ceiling varies roughly 2$\times$ with working-set size. Every ``\% of peak'' we
had published for a cache-resident shape was divided by the DRAM ceiling and was
therefore inflated (Sec.~\ref{sec:ceiling}).

\item \textbf{Direct measurement of streaming-mode overhead}, replacing a figure
that had been inferred from a curve fit and had silently absorbed an unrelated
bug (Sec.~\ref{sec:crossover}).

\item \textbf{Runtime code generation verified by encoding diff} against the
linked hand-written kernels, and integration as a primitive in a tiled-execution
IR runtime (Sec.~\ref{sec:jit}).

\item \textbf{A correctness-gated external comparison} against PyTorch and
ExecuTorch, in which the interesting finding concerns fusion rather than kernel
quality (Sec.~\ref{sec:external}).
\end{enumerate}

% =========================================================================
\section{Background and Target}

\subsection{SME and Streaming SVE on the M4}

SME adds a streaming execution mode (\code{PSTATE.SM}) with a streaming vector
length (SVL) that is independent of any non-streaming vector length, plus a
two-dimensional accumulator array, ZA. Our target reports
\code{hw.optional.arm.FEAT\_SME} and \code{hw.optional.arm.FEAT\_SME2}, and
SVL${}=512$ bits, i.e.\ 16 FP32 lanes, queried at run time with \code{RDSVL}.

Three properties of this implementation shape every kernel below.

\emph{No gather/scatter in streaming mode.} Vector-offset loads and stores
require \code{FEAT\_SME\_FA64}, which Apple does not implement; using them
inside a streaming region raises \code{SIGILL}. A strided row therefore cannot
be loaded in one instruction.

\emph{Streaming entry destroys register state.} \code{SMSTART}/\code{SMSTOP}
zero the Z registers and hence \code{d8}--\code{d15}, which AAPCS64 makes
callee-saved. Any value a kernel needs across the transition must live in
memory, and a kernel that does not save all eight registers is
ABI-nonconforming even when a standalone benchmark appears to work.

\emph{The feature set must be detected, not assumed.} Our own project documents
asserted ``the M4 is SME1'' for several sprints while the hardware reported
\code{FEAT\_SME2}${}=1$. Every benchmark run now prints a provenance header
containing the \code{sysctl} values, the git commit (marked \code{-dirty} on a
modified tree), compiler and build type.

\subsection{The two norms, stated precisely}

Over a feature axis of size $N$, with $\varepsilon$ a small constant:
\begin{align}
\text{LayerNorm:}\quad & y = \gamma \odot \frac{x-\mu}{\sqrt{\sigma^2+\varepsilon}} + \beta \\
\text{RMSNorm:}\quad   & y = \gamma \odot \frac{x}{\sqrt{\tfrac{1}{N}\sum_i x_i^2 + \varepsilon}}
\end{align}

The distinction that matters for bandwidth is \emph{reduction stages} versus
\emph{input traversals}; the common shorthand ``two-pass / single-pass''
conflates them. LayerNorm needs two reduction stages (mean, then variance from
\emph{centred} values) plus output generation, hence three traversals of the
input and 3R+1W of traffic. RMSNorm needs one reduction stage plus output,
hence two traversals and 2R+1W (Fig.~\ref{fig:traversals}). That $4{:}3$ ratio
is the structural source of
RMSNorm's advantage and is why the 2R+1W model must not be applied to
LayerNorm.

\begin{figure}[t]
\centering
\includegraphics[width=\columnwidth]{figures/traversals}
\caption{Reduction stages versus input traversals. LayerNorm needs two
reduction stages and therefore three traversals; RMSNorm needs one and therefore
two. The resulting $4{:}3$ traffic ratio is structural.}
\label{fig:traversals}
\end{figure}

\begin{table}[t]
\caption{Per element of output. Operational intensity is computed against moved
bytes. Reaching the M4's $\approx$1984\,GFLOPS FMOPA peak would require
$\approx$31\,flop/B, about 100$\times$ more than either norm offers.}
\label{tab:intensity}
\centering
\begin{tabular}{lcc}
\toprule
 & RMSNorm & LayerNorm \\
\midrule
reduction stages          & 1 ($\Sigma x^2$) & 2 (mean, variance) \\
input traversals          & 2                & 3 \\
traffic, moved            & 2R+1W = 12\,B    & 3R+1W = 16\,B \\
traffic, useful (1R+1W)   & 8\,B             & 8\,B \\
FP operations             & $\approx$3       & $\approx$5 \\
operational intensity     & $\approx$0.25\,flop/B & $\approx$0.31\,flop/B \\
\bottomrule
\end{tabular}
\end{table}

We keep three claims separate throughout, because they are separate:
\emph{semantics} (RMSNorm omits mean-centering, hence one fewer traversal);
\emph{literature} (Zhang and Sennrich~\cite{zhang2019rmsnorm} report comparable
task performance with runtime reductions of roughly 7--64\,\% across their
experiments); and \emph{our measurement} (1.8--2.3$\times$ on the shapes
evaluated here). Kernel numerical agreement is \emph{not} model accuracy: our
tests show that our RMSNorm matches an \emph{RMSNorm} reference, which says
nothing about whether substituting one norm for the other preserves accuracy in
a network.

% =========================================================================
\section{Kernel Design}

\subsection{Layout and loop order}

Activations are stored column-major with an explicit leading dimension,
\code{a[row + col*ld]}, following the surrounding lab code. A column is
therefore contiguous and a row is strided.

The natural loop order---outer over rows, inner reduction over the feature
axis---would require gathering a strided row, which streaming mode cannot do
(Sec.~II-A). We therefore transpose the loop order: the outer loop walks
\emph{blocks of SVL rows} and the inner loop walks all $N$ columns. Each inner
load is one contiguous \code{LD1W} covering SVL consecutive rows of a single
column, and $\gamma$, $\beta$ enter through \code{LD1RW} broadcasts. This is the
structural decision from which everything else follows: the reduction becomes
\emph{vertical} across lanes, so each lane accumulates one row independently and
no horizontal reduction instruction is needed at all.

\subsection{Vector-length agnosticism}

No kernel hard-codes SVL. \code{CNTW} supplies the lane count, \code{WHILELO}
builds the row predicate so that full blocks and the final partial block share
one code path, and \code{ADDVL}/\code{INCW} advance pointers and indices by one
vector length. The row granularity of our winning variants is $4\times$SVL---64
rows on this machine---but it is derived at run time, never written as a
literal. The C++ harness obtains the same quantity through \code{RDSVL} rather
than assuming~16.

\subsection{One streaming region per call}

The entire norm executes inside a single \code{SMSTART}/\code{SMSTOP} pair.
Section~\ref{sec:crossover} quantifies what that decision is worth.

\subsection{Precision and the reference}

Kernels accumulate in FP32, the precision of the data. The reference
implementation accumulates in double and is kept deliberately unoptimized: it is
the correctness oracle, not a performance target. LayerNorm computes variance
from centred values, never as $E[x^2]-\mu^2$; Sec.~\ref{sec:numerics} shows what
that choice buys.

% =========================================================================
\section{Measurement Methodology}
\label{sec:method}

Three rules, each adopted after a measurement error forced it.

\subsection{Correctness gates performance, per configuration}

No throughput figure is produced for code that has not first matched the float64
reference \emph{on that exact shape}. The gate runs the function that is about
to be timed, compares its complete output, and refuses to produce a timing on
disagreement. An earlier version verified each row but still printed GiB/s
beside a failure flag, which a reader scanning the throughput column would never
notice. The submitted run reports \textbf{66/66 configurations verified before
timing}, and the external baselines 6/6 each.

\subsection{Bytes are counted one way}

All GiB/s figures count \emph{useful bytes}: one read and one write per element,
the algorithm's minimum traffic, for kernels and bandwidth probes alike. Moved
bytes are $1.5\times$ that for RMSNorm and $2.0\times$ for LayerNorm. Keeping
the useful-bytes convention makes the extra traversal visible as a lower
percentage of peak, which is exactly the gap the access-density lever attacks,
rather than a number to hide.

\subsection{The ceiling is a curve, not a constant}
\label{sec:ceiling}

This is the project's largest correction and it invalidated percentages across
five report sections.

Our roofline probe is a STREAM-style~\cite{mccalpin1995stream} scale-add
executed the way the kernels access memory: one streaming region, contiguous \code{LD1W}/\code{ST1W},
predicated tail. We originally measured it at a single 128\,MiB-per-array
footprint, obtaining 58.8\,GiB/s, and then divided \emph{every} shape by that
number.

That figure is the \emph{DRAM} ceiling. Sweeping the same probe across
working-set sizes (Fig.~\ref{fig:ceiling}) shows it reaching 115--117\,GiB/s
between 512\,KiB and 16\,MiB. A cache-resident kernel never faces DRAM
bandwidth, so dividing it by the DRAM ceiling credits it for a constraint it
never met. The most-quoted number in our earlier drafts---RMSNorm V6 at
``$\approx$95\,\% of the moved-bytes roofline''---is really $\approx$50\,\%.

The methodological point on which the correction rests: at a small footprint
both arrays stay resident after the first repetition, so the probe measures
\emph{steady-state resident} bandwidth. That is a legitimate denominator only
because the norm kernels are benchmarked identically, with repeated calls over
the same buffers. Probe and kernel must sit in the same residency regime or
their ratio is meaningless. The harness now sweeps the ceiling and each table
divides by the ceiling measured at its own working set, so the numbers
regenerate from the tree instead of depending on a constant that suits one
regime.

Two further readings fall out of the sweep. Streaming mode costs $\approx$25\,\%
against NEON at DRAM but is nearly free in cache ($0.93\times$), so ``streaming
is a structural handicap'' is a DRAM-regime statement rather than a property of
the mode. And below $\approx$512\,KiB the SSVE figure stops being a bandwidth
number at all: NEON climbs toward L1 speed while SSVE falls, because
\code{SMSTART} becomes a visible share of a sub-microsecond pass.

This is the same class of error we had ourselves named earlier in the project
(``peak of \emph{what}?''), when we established that a NEON probe is the wrong
ceiling for a streaming-mode kernel. That fixed the execution-mode half of the
question and left the footprint half unasked.

\begin{figure}[t]
\centering
\includegraphics[width=\columnwidth]{figures/ceiling_curve}
\caption{The bandwidth ceiling is a curve, not a constant. Apple M4, single
core. Until this was measured, every ``\% of peak'' divided by the dotted line,
which is the ceiling at a 256\,MiB footprint. The gap between the dotted line
and the dashed curve over the cache-resident range is the $\approx$2$\times$
error in our own published efficiency figures.}
\label{fig:ceiling}
\end{figure}

\subsection{Reporting}

Headline figures are the best of $N$ repetitions, an estimate of what the kernel
does when nothing interferes. Because a minimum alone cannot distinguish a
typical result from a lucky outlier, we report the median and the p10--p90 band
from the same samples beside it. They agree to about 1\,\%, so the best-case
figures this project has reported throughout \emph{were} typical---previously an
assumption, now shown. Repetition counts scale with working-set size rather than
being fixed: at a fixed count the smallest shape once reported a hand-written
kernel 10\,\% below its own word-identical JIT twin, which is impossible for
identical code and was therefore sampling noise.

% =========================================================================
\section{Optimization and Ablation}
\label{sec:ablation}

Both norms were taken through the same ladder, one lever per variant, each rung
verified and measured on the same harness and the same shapes.

\begin{table}[t]
\caption{Consolidated ablation, useful bytes (1R+1W), best-of-$N$. Percentages
are against the ceiling measured \emph{at each shape's own footprint}
(115.60 and 58.78\,GiB/s respectively). Every row was verified against the
float64 reference before it was timed. This table is the consolidated run, in
which each rung is measured once in sequence; the per-lever percentages quoted
in the text come from the dedicated order-alternating A/B studies, which is the
instrument that controls for drift in a paired comparison.}
\label{tab:ablation}
\centering
\small
\begin{tabular}{lrrrr}
\toprule
& \multicolumn{2}{c}{RMSNorm} & \multicolumn{2}{c}{LayerNorm} \\
\cmidrule(lr){2-3}\cmidrule(lr){4-5}
stage & 16\,MiB & 256\,MiB & 16\,MiB & 256\,MiB \\
\midrule
scalar reference & 0.90 & 0.59 & 0.84 & 0.47 \\
V0 SSVE & 24.87 & 10.04 & 12.82 & 7.38 \\
V1 \code{FRSQRTE}+NR & 24.88 & 10.06 & 12.79 & 7.38 \\
V2 hoisted $1/N$ & 24.82 & 10.29 & 14.40 & 7.39 \\
V3 $2\times$ unroll & 24.84 & 10.34 & --- & --- \\
V4 accumulator ILP & 31.37 & 10.67 & 16.36 & 7.69 \\
V5 load pipelining & 31.35 & 10.65 & 16.35 & 7.56 \\
\textbf{V6} access density & \textbf{38.46} & \textbf{21.05} & \textbf{16.47} & \textbf{13.09} \\
\textbf{V7} SME2 multi-vector & \textbf{38.45} & \textbf{24.63} & \textbf{16.47} & \textbf{13.52} \\
Welford (single stage) & --- & --- & 12.07 & 6.67 \\
ZA residency & 24.86 & 10.29 & 14.41 & 7.39 \\
JIT (automatic ISA) & 38.46 & 24.68 & 16.48 & 13.49 \\
\midrule
\% of footprint ceiling (V7) & 33.3 & 41.9 & 14.2 & 23.0 \\
\bottomrule
\end{tabular}
\end{table}

\subsection{What each lever bought}

\textbf{V1, reciprocal-square-root arithmetic} ($\approx$0). Replacing
\code{FSQRT}+\code{FDIV} with \code{FRSQRTE} plus one Newton--Raphson step
halves the latency of the inverse but changes nothing measurable: it is off the
critical path.

\textbf{V2, hoisting $1/N$} ($\approx$0 for RMSNorm, \textbf{+9--12\,\% for
LayerNorm}). This is the one verdict that does \emph{not} transfer between the
norms, and the asymmetry is structural. RMSNorm performs one division per row
block, already amortized over the column loop. LayerNorm performs two, and both
sit at serialization points \emph{between} its three passes---the mean gates
pass 2, the variance gates pass 3---so they are a far larger share of runtime.

\textbf{V3, unrolling} ($\approx$0). The out-of-order core already overlaps
sequential column loads.

\textbf{V4, accumulator ILP} (\textbf{+27--33\,\%}). Four independent
accumulator chains replace one. This is the first real win, and it retroactively
explains V3: unrolling without breaking the dependency chain cannot help,
because the chain was the bottleneck. The four partial sums combine in tree
order rather than the reference's sequential order, so V4 and V5 are verified at
a documented wider tolerance rather than the strict one---an honest statement of
the accuracy they meet, not a gate quietly widened.

\textbf{V5, load pipelining} ($\approx$0 versus V4, as pre-registered). Register
renaming already hoists the loads once the chains are broken. This closes the
``more outstanding loads'' family; explicit prefetch works the same lever one
step earlier and receives a reasoned skip.

\textbf{V6, access density} (\textbf{+131\,\% RMSNorm, +71\,\% LayerNorm} versus
V0 in the DRAM regime). This is the winning lever, and it is \emph{not} the one
we planned. We had intended a residency optimization---keep the row on core so
the second traversal reads from cache---and measured the premise first. Pass-2
residency was \emph{already} satisfied: the reuse distance between traversals of
one row block is only $64N$ bytes, L2-resident by construction. What actually
collapses throughput is the \emph{density} of the column walk. Each column touch
reads 64\,B out of every $4M$ bytes, and at an equal 32\,MB footprint the aspect
ratio alone moves throughput from 31 to 18\,GiB/s.

V6 therefore processes four consecutive SVL-row blocks per outer iteration, so
each column touch becomes four back-to-back loads at consecutive
\code{MUL VL} offsets---256 contiguous bytes, four times denser. One accumulator
\emph{per row block} preserves V4's four independent chains while restoring the
reference's sequential summation order, so V6 needs no reassociation tolerance.

\textbf{V7, SME2 multi-vector accesses} (\textbf{+17.2\,\% RMSNorm at DRAM},
+2.7\,\% LayerNorm, $\approx$0 cache-resident). V6's group loop already touches
four consecutive row blocks per column, which is exactly the operand shape of
SME2's four-vector contiguous access, so V7 is V6 with four loads folded into
one and four stores into one. Same addresses, same traffic, same arithmetic,
same summation order: the only variable is instruction count, which makes it a
clean test of whether V6 is issue-bound or memory-bound. Because no value
changes, the correctness gate is \emph{bit-identity} with V6, not a tolerance.

Folding the accesses also eliminates the separate tail path. V6 needs one
because its group path is unpredicated, and that tail touches 64\,B per column
instead of 256\,B---precisely the low density V6 exists to fix. SME2's
predicate-as-counter (\code{WHILELO}\,\dots\,\code{VLx4}) describes
$\min(M-i,\,4\cdot\mathrm{SVL})$ active elements \emph{across} the group, so one
loop body serves full and final partial groups alike.

\begin{figure}[t]
\centering
\small
\begin{verbatim}
V6 -- four accesses per column touch:
  ld1w {z0.s},  p0/z, [x8]
  ld1w {z24.s}, p0/z, [x8, #1, mul vl]
  ld1w {z25.s}, p0/z, [x8, #2, mul vl]
  ld1w {z26.s}, p0/z, [x8, #3, mul vl]
  fmla z2.s,  p0/m, z0.s,  z0.s
  fmla z16.s, p0/m, z24.s, z24.s
  fmla z17.s, p0/m, z25.s, z25.s
  fmla z18.s, p0/m, z26.s, z26.s

V7 -- the same 256 B in one instruction:
  ld1w {z0.s - z3.s}, pn8/z, [x8]
  fmla z4.s, p0/m, z0.s, z0.s
  fmla z5.s, p0/m, z1.s, z1.s
  fmla z6.s, p0/m, z2.s, z2.s
  fmla z7.s, p0/m, z3.s, z3.s
\end{verbatim}
\caption{The V6$\rightarrow$V7 transformation, reduction loop. Same addresses,
same traffic, same arithmetic, same summation order; only the load-instruction
count changes. Because no value changes, the correctness gate is bit-identity
with V6 rather than a tolerance. The counter predicate \code{pn8}, built by
\code{WHILELO}\,\dots\,\code{VLx4}, also governs the final partial group, which
is what removes V6's separate tail path entirely.}
\label{fig:v7}
\end{figure}

\begin{figure}[t]
\centering
\includegraphics[width=\columnwidth]{figures/ablation_dram}
\caption{The ablation in the true-DRAM regime (256\,MiB), where the levers
separate most clearly. Access density (V6) is what moves the needle; SME2
multi-vector accesses (V7) add to it. ZA residency lands back at V0 level
despite moving strictly fewer bytes.}
\label{fig:ablation}
\end{figure}

\subsection{The ZA tile: a measured, explained negative}
\label{sec:za}

SME's ZA array offers roughly 4\,KB of on-core storage beyond the 32 Z
registers---enough to hold one SVL-row block of activations. Staging $x$ there
during the reduction and reading it back afterwards converts RMSNorm's 2R+1W
into 1R+1W, and LayerNorm's 3R+1W into 1R+1W. The traffic saving is real.

Both lose, decisively: 39--65\,\% slower than V6 for RMSNorm and 51--68\,\% for
LayerNorm. The mechanism is that each element now costs two (RMSNorm) or three
(LayerNorm) additional \code{mova} operations, and \code{mova} throughput, not
memory, is what binds at these shapes. The clinching observation is internal to
the kernel: its own streaming \emph{fallback} path outruns its ZA-resident fast
path. We also verified that no footprint rescues it---at 32--64\,MB the ZA path
stays pinned near 10\,GiB/s while V6 delivers 22--25---because V6's grouping
keeps the second traversal cache-resident even there. The DRAM read that ZA
``saves'' was never a DRAM read.

Two aspects of this result are worth more than the numbers.

First, that LayerNorm loses by a \emph{larger} margin than RMSNorm was
pre-registered before measuring: a 50\,\% traffic cut that costs three
\code{mova} per element rather than two must clear a proportionally larger tax.
The prediction was recorded, then confirmed, which is what makes it evidence
rather than rationalization.

Second, we later \emph{retracted our own reasoning}. We had declined to try SME2
multi-vector \code{mova} on the argument that ``even a 2$\times$ speed-up only
reaches a tie''. That rested on an untested factor: \code{mova} had never been
characterized as issue-bound versus ZA-port-bound. Measured directly, the
single-vector form reaches 3.88\,G\,vec/s and the four-vector form
15.50---exactly $4.00\times$. It is issue-bound and the fold is $4{:}1$, so our
own arithmetic did not hold. We built both kernels: they gain \textbf{+62\,\%}
(RMSNorm) and \textbf{+114\,\%} (LayerNorm) and \emph{still} lose to V7 by
37\,\% and 24\,\%. The verdict is unchanged but is now a measurement rather than
an extrapolation, and the residual cause differs from our original diagnosis:
the ZA kernels touch one row block per iteration (64\,B per column) against
V6/V7's four (256\,B), so what limits them is the same access-density lever, not
\code{mova}.

\subsection{V7's mechanism was falsified}

Two hypotheses explain V7's DRAM win and predict opposite LayerNorm outcomes.
``Fewer instructions retire faster'' predicts a \emph{larger} gain for
LayerNorm, which folds more accesses per element. ``Relieved memory-level
parallelism''---one instruction requesting 256 contiguous bytes occupying less
load-tracking capacity than four requesting 64---predicts a \emph{smaller} one,
since LayerNorm does more FP work per byte and is less memory-starved.
LayerNorm gained 2.7\,\% against RMSNorm's 17.2\,\%, which falsifies the
instruction-count explanation.

We then built a two-vector control: identical work, identical 256\,B per column,
issued as two 128\,B accesses instead of one 256\,B access, so
bytes-per-instruction is the only variable. A proportional
memory-level-parallelism model predicts it lands halfway. It does not
(Table~\ref{tab:v7x2}): the curve \emph{saturates}, with $4\!\to\!2$ capturing
90\,\% of the win and $2\!\to\!1$ adding 1.5\,\%. That is inconsistent with our
own pre-registered model \emph{and} with a ``256\,B burst matches the
prefetcher'' model. Both explanations are dead. What survives is narrower than
either: a threshold effect on load-instruction count, in load-dominated loops,
in the latency-exposed regime only. Why the threshold sits between four and two
is not established here and would need microarchitectural counters this harness
does not read.

\begin{table}[t]
\caption{The two-vector control. RMSNorm, 256\,MiB. The gain saturates, which
falsifies both pre-registered explanations.}
\label{tab:v7x2}
\centering
\begin{tabular}{lccr}
\toprule
variant & load insns/column & GiB/s & vs V6 \\
\midrule
V6   & 4 ($4\times$64\,B)  & 20.94 & --- \\
V7x2 & 2 ($2\times$128\,B) & 24.31 & +16.1\,\% \\
V7   & 1 ($1\times$256\,B) & 24.67 & +17.9\,\% \\
\bottomrule
\end{tabular}
\end{table}

\subsection{What binds, and where}

Combining the ablation with issue-rate measurements gives a regime-dependent
answer rather than a single one. We measured the SSVE floating-point issue
ceiling at 0.97\,G vector instructions/s (equivalently 31.0\,GFLOPS for
\code{FMLA}), and \code{FMUL} issues at exactly the same rate. At 16\,MiB,
RMSNorm~V7 sustains \emph{100\,\%} of that ceiling; at 256\,MiB it reaches
67\,\%. Cache-resident, the kernel is FP-issue-bound, not memory-bound---which
is precisely why V7's multi-vector loads buy nothing there and 17\,\% at DRAM.

This is why we describe these kernels as \emph{bandwidth-dominated for the
shapes measured} rather than flatly memory-bound. Against footprint-matched
ceilings the best kernel reaches $\approx$33\,\% useful / $\approx$50\,\% moved,
not $\approx$95\,\%, and reduction dependency chains, FP issue rate,
load-instruction count and reciprocal-square-root latency all appear in the
ablation.

% =========================================================================
\section{Numerical Behaviour}
\label{sec:numerics}

The project asserted from the outset that single-pass variance is dangerous and
that LayerNorm's centred formulation avoids it---but never implemented the
dangerous version, so the claim had no counterexample. We supplied one: the
naive $E[x^2]-\mu^2$, the centred two-pass, and Welford's recurrence, all
accumulating in FP32 (the kernels' precision; a float64 demonstrator would prove
nothing about them) against a float64 oracle. The stress axis is a shift, which
leaves the variance untouched but scales the condition number
$\kappa = 1 + \mu^2/\sigma^2$ by roughly $s^2$, isolating conditioning from
every other property of the data. Summation order is trustworthy because the
build sets no fast-math flags, so IEEE semantics forbid reassociating the
reductions.

\begin{table}[t]
\caption{FP32 relative error against a float64 oracle, $N=512$, identical input.
Past shift $10^5$ the naive estimator returns a \emph{negative} variance, so a
kernel built on it emits \code{NaN}.}
\label{tab:stability}
\centering
\begin{tabular}{lrrr}
\toprule
shift & naive 1-stage & two-pass & Welford \\
\midrule
$0$      & 5.09e--07 & 3.89e--07 & 2.07e--07 \\
$10^{1}$ & 1.39e--04 & 3.98e--07 & 1.98e--07 \\
$10^{2}$ & 8.52e--03 & 1.18e--07 & 4.16e--07 \\
$10^{3}$ & 1.25e+00 & 3.75e--07 & 4.73e--06 \\
$10^{4}$ & 1.59e+02 & 2.81e--05 & 5.41e--05 \\
$10^{5}$ & 1.02e+04$^{\dagger}$ & 1.51e--07 & 8.94e--05 \\
$10^{6}$ & 6.80e+06$^{\dagger}$ & 7.43e--07 & 6.71e--03 \\
\bottomrule
\multicolumn{4}{l}{\footnotesize $^{\dagger}$ negative variance; $\sqrt{\cdot}$ is \code{NaN}.}
\end{tabular}
\end{table}

Three results follow. \emph{Naive error tracks $\kappa$}: each row multiplies
the shift by ten and hence $\kappa$ by about a hundred, and the naive column
rises by the same factor---the classical conditioning result reproduced on this
hardware. \emph{The failure is qualitative, not gradual}: from shift $10^5$ the
estimate goes negative, and a negative variance has no square root, so the
kernel produces \code{NaN} rather than an inaccurate number. That is why we
treat stability as part of correctness rather than as a separate concern.
\emph{Centred two-pass is flat across twelve orders of magnitude of $\kappa$},
and that immunity is exactly what LayerNorm's third traversal buys---accuracy
bought with bandwidth, priced.

\begin{figure}[t]
\centering
\includegraphics[width=\columnwidth]{figures/stability}
\caption{Where each FP32 variance formulation breaks, $N=512$. Above the dotted
line the error exceeds the quantity being estimated. The naive estimator's error
tracks the condition number and returns a \emph{negative} variance from shift
$10^5$; the centred two-pass is flat across twelve orders of magnitude, and that
immunity is what LayerNorm's third traversal buys.}
\label{fig:stability}
\end{figure}

\begin{table}[t]
\caption{The shipped kernels on the same stress data, $M=64$, $N=512$;
max$|$error$|$ against the float64-accumulating reference. V0 uses exact
\code{FSQRT}, V6 uses \code{FRSQRTE}+NR.}
\label{tab:kernelacc}
\centering
\small
\begin{tabular}{lrrrr}
\toprule
shift & LN V0 & LN V6 & LN Welford & RMS V6 \\
\midrule
$0$      & 7.15e--07 & 1.72e--05 & 1.72e--05 & 1.73e--05 \\
$10^{2}$ & 1.30e--04 & 1.30e--04 & 1.02e--04 & 6.79e--06 \\
$10^{4}$ & 5.37e--03 & 5.38e--03 & 2.16e--02 & 3.10e--06 \\
$10^{5}$ & 1.23e--02 & 1.23e--02 & 2.14e--01 & 5.25e--06 \\
$10^{6}$ & 2.07e--03 & 2.08e--03 & 3.02e--01 & 5.19e--06 \\
\bottomrule
\end{tabular}
\end{table}

Scalar estimators prove nothing about SSVE kernels, so we placed the shipped
kernels on the same axis (Table~\ref{tab:kernelacc}). Three findings refine the
story. The
\code{FRSQRTE}+NR substitution costs accuracy---about $24\times$ at shift
0---but only where the problem is already well conditioned; from shift $10^2$
upward the exact-\code{FSQRT} and approximate variants are indistinguishable.
What limits LayerNorm on hard data is therefore \emph{not the variance at all}:
the residual error lives in the output expression $(x-\mu)$, where FP32 carries
the difference of two large numbers in only a handful of bits. At shift $10^5$
the predicted ULP is 1.6e--2 relative to a unit signal and the measured error is
1.2e--2; that the exact and approximate kernels coincide there is the evidence
that the square-root arithmetic is not the binding constraint. Finally,
\textbf{RMSNorm is flat at $\approx$5e--6 across the whole sweep}, about
$2300\times$ better than LayerNorm at shift $10^5$, because it never forms a
mean and so never performs the cancelling subtraction---in the reduction or in
the output. Its stability is structural, and it is the same property that makes
it faster. Accuracy and throughput point the same way here, which is unusual
enough to be worth stating.

The honest counterpart: RMSNorm has its own limit. $\sum x^2$ must fit in FP32,
so the reduction overflows once $N x^2 > \textsc{flt\_max}$, i.e.\ at
$|x|\approx\sqrt{\textsc{flt\_max}/N}$. A test straddles that boundary at
$\pm$10\,\% to show the failure is where the arithmetic says it is.

We also corrected a claim of our own here. We had reported Welford as
``$\approx$100$\times$ less accurate'' than two-pass. Re-measured across more
shift magnitudes, that holds at exactly one shift (and the factor is
$\approx$46$\times$), is a wash at two others, and \emph{reverses} at $10^5$
where Welford is $2.5\times$ better. The stated mechanism was wrong too:
Welford's \emph{variance} is the more accurate of the two, and its extra output
error comes from the running \emph{mean}, which updates at full input magnitude
every element. What does not change is the decision---two-pass remains the
incumbent because Welford is 25--30\,\% slower for a measured reason, a
per-column scalar division that serializes the loop, despite moving 25\,\% fewer
bytes. That is a speed argument and it is unaffected.

% =========================================================================
\section{Code Generation and Integration}
\label{sec:jit}

The kernels are generated at run time and invoked through a tiled-execution-IR
(TEIR) runtime, which is the point of the exercise: the kernel is a leaf the
compiler schedules, not a standalone program.

\subsection{Verification by encoding diff}

A generator can emit correct instruction words that its author does not
understand, and ``it assembles'' is not ``it is the instruction you meant''.
Each emitter is a one-to-one transcription of a hand-written kernel---same
registers, same order---and the generated buffer is compared word by word
against the \emph{linked} hand-written kernel, read from its function address at
run time so the reference can never go stale. The diff is green at 157/157 and
208/208 words for the SME1 variants and 110/144 for the SME2 variants, so the
generated kernel inherits the verification of one that already passed the full
suite instead of re-earning it.

Building the encoders this way---golden words extracted from the toolchain
first, every encoder pinned to one---also surfaced two latent defects in
pre-existing lab code that behavioural tests had passed over for weeks: a
predicate constructor that emitted a byte-element predicate where it documented
a word-element one (functionally masked, since an all-true predicate of either
kind governs the same operations), and a streaming-mode entry helper that
enabled ZA as well as SM contrary to its name. An encoding diff checks what code
\emph{is}; behavioural tests check what it \emph{does}, and the first catches a
class of error the second structurally cannot.

\subsection{Feature-dependent emission}

The generator takes an ISA parameter. With \code{automatic}---what a real caller
gets---it emits the SME2 variant where \code{FEAT\_SME2} is present and the SME1
variant otherwise; explicit values let the tests pin either path regardless of
build host. Emission costs $\approx$4\,$\mu$s once, against per-call execution
times from 0.4 to 135\,$\mu$s, and the returned pointer is reused.

\subsection{Loop-nest integration and threading}

The norms are registered as TEIR primitives that the runtime places in a
generated loop nest, honouring strides and leading dimensions. Marking the row
axis parallel gives thread scaling without touching the kernel, which remains
single-tile (Table~\ref{tab:threading}).

Two integration findings are worth recording. First, this is where the
\code{d8}--\code{d15} ABI violation stopped being latent: our benchmark harness
had masked it caller-side with \code{volatile}, and a generated loop nest has no
such workaround. Fixing it properly---saving the full callee-saved range around
the streaming region, in both the assembly and the generator, then re-running
the encoding diff to confirm the two stay bit-identical---was a prerequisite for
integration rather than cleanup.

Second, work distribution interacts with kernel granularity. A naive
$M/\text{threads}$ split hands every thread a partial group whenever the
quotient is not a multiple of the 64-row granularity, and those thread counts
measured $\approx$45\,\% slower than their neighbours. We first attributed this
to core heterogeneity; isolating the effects showed partial groups cost nothing
by themselves, and the cliff persisted after deleting the tail path entirely.
Distributing in whole row-blocks removes it empirically---+39\,\% at six threads
and +42\,\% at ten, and 0\,\% at thread counts that were already aligned, which
is the control that rules out a generic ``aligned is faster'' explanation. We
record that the underlying mechanism remains \emph{unidentified}: a
cache-line-aligned base offset is still 11\,\% slower than the allocation start,
so it is not simply line alignment. The fix is justified by its control, not by
our explanation.

\begin{table}[t]
\caption{TEIR-invoked, OpenMP row-parallel, 256\,MiB working set. Every
configuration verified against the reference before its figure is reported. The
chip-wide ceiling is 84.30\,GiB/s.}
\label{tab:threading}
\centering
\begin{tabular}{lrrrr}
\toprule
& \multicolumn{2}{c}{RMSNorm} & \multicolumn{2}{c}{LayerNorm} \\
\cmidrule(lr){2-3}\cmidrule(lr){4-5}
threads & GiB/s & \% chip & GiB/s & \% chip \\
\midrule
1  & 24.67 & 29.2 & 13.48 & 15.9 \\
2  & 29.36 & 34.8 & 18.02 & 21.4 \\
4  & 35.86 & 42.5 & 22.62 & 26.8 \\
8  & 40.98 & 48.6 & 24.86 & 29.5 \\
16 & 43.57 & 51.6 & 26.07 & 30.9 \\
\bottomrule
\end{tabular}
\end{table}

Scaling saturates well before the ceiling---$1.76\times$ and $1.93\times$ at 16
threads. This is expected rather than suspicious: an independent
characterization of the same machine finds that a single performance core can
largely saturate the cluster's SME resource, so the M4 is not well modelled as
ten cores times single-core throughput.

% =========================================================================
\section{When the Kernel Is Worth Using}
\label{sec:crossover}

Every per-call cost this project had quoted was \emph{inferred}, from the
intercept of a linear fit. That instrument had already failed once: we read
$t_0\approx1\,\mu$s from it, and later found that about 70\,\% of that was an
uncached feature-detection syscall in a wrapper, not streaming cost at all. The
fit had silently absorbed an unrelated bug.

We therefore measured the transition directly, against an identical
transition-free loop as a control so that loop overhead is subtracted rather
than attributed. One round trip through streaming mode costs \textbf{9.04\,ns},
and the SM+ZA form is indistinguishable from SM-only---so managing the ZA state
is free, and the ZA kernels' problem was never the transition. The per-call
floor, measured rather than fitted at the smallest possible call and timed in
batches because a single call sits at clock resolution, is 45.8\,ns in the
submitted run, of which the transition is 19.7\,\%.

A linear fit is deliberately \emph{not} used here: over these sizes it returns a
negative intercept, which is the model announcing that it does not apply.

\begin{figure}[t]
\centering
\includegraphics[width=\columnwidth]{figures/crossover}
\caption{Crossover against the scalar reference at $N=512$. The SME kernel's
time is flat until the group is full, which is what sets the crossover.}
\label{fig:crossover}
\end{figure}

\begin{table}[t]
\caption{RMSNorm at $N=512$, sweeping rows. The SME kernel's time is flat from
$M=1$ to $M=64$ because that is one group; at $M=256$ it is exactly
$4\times$ the $M=64$ time.}
\label{tab:crossover}
\centering
\begin{tabular}{rrrrr}
\toprule
$M$ & scalar ($\mu$s) & V7 ($\mu$s) & speedup & groups \\
\midrule
1   & 0.250   & 9.291  & 0.03$\times$ & 1 \\
4   & 2.416   & 6.291  & 0.38$\times$ & 1 \\
8   & 4.916   & 6.209  & 0.79$\times$ & 1 \\
\textbf{16}  & 9.916 & 6.250 & \textbf{1.59$\times$} & 1 \\
64  & 74.833  & 6.166  & 12.14$\times$ & 1 \\
256 & 422.208 & 25.250 & 16.72$\times$ & 4 \\
\bottomrule
\end{tabular}
\end{table}

The crossover against the scalar reference is at $M\approx16$ rows
(Fig.~\ref{fig:crossover}), and the
cause is \emph{not} streaming overhead. A 9\,ns transition cannot explain a
6\,$\mu$s plateau; group granularity explains it exactly. Our winning variants
process four SVL-row blocks at a time---64 rows here---and a partially filled
group costs the same as a full one, because unused lanes are predicated off
rather than skipped. At $M=1$ the kernel performs 64 rows of work to produce one
row of result.

This also quantifies a design decision previously taken on intuition. The
alternative to one streaming region per call is one per row, which at $M=128$
costs $128\times9.04\,\text{ns}\approx1.16\,\mu$s---about twenty-five times the
entire per-call floor. The decision is right, and now right by a measured
margin. But the constant is 9\,ns: \emph{``streaming mode is expensive'' is true
per row and false per call}, and only the second is what these kernels do.

% =========================================================================
\section{External Comparison}
\label{sec:external}

\subsection{Method, and why it is this strict}

An earlier vendor comparison in this project produced four wrong claims at once:
a ``vendor RMSNorm'' that was our own composition of two primitives, an
equivalence between a library routine and LayerNorm that ignored $\varepsilon$,
$\gamma$ and $\beta$, and a headline speed-up that measured a \emph{layout
mismatch}---the library forced to walk our column-major matrix with a
cache miss per element. Given its own contiguous layout, that library was in
fact faster than us. Every rule below exists because of one of those failures.

\emph{Native layout for everyone.} Our kernels are column-major over a strided
feature axis; the libraries normalize the last, contiguous axis of a row-major
tensor. Both are efficient in their own layout, so each is measured in its own,
and the transpose lives only in the correctness checker, never in a timed path.
This answers \emph{what throughput each implementation achieves given the layout
it was designed for}, and explicitly not \emph{what it would cost to substitute
our kernel for a framework operator on the framework's own representation},
which would have to include boundary conversion.

\emph{Correctness gates the comparison.} Every library output is dumped and
verified against our float64 reference before its throughput is quoted; all
agree to $\le$3.1e--6, so the numbers compare implementations of the same
function. Both drivers dump every benchmarked shape and write a manifest that
the checker walks, so the set of shapes timed and the set gated cannot drift
apart.

\emph{One thread, same shapes, same byte convention, pinned versions.}

\begin{figure}[t]
\centering
\includegraphics[width=\columnwidth]{figures/baselines}
\caption{Against general-purpose frameworks at the 256\,MiB shape,
single-threaded, each implementation in its native layout.}
\label{fig:baselines}
\end{figure}

\begin{table}[t]
\caption{Single-threaded, useful bytes, each implementation in its native
layout, every shape verified to compute the same function before timing.
PyTorch~\cite{pytorch} 2.13.0, ExecuTorch~\cite{executorch} 1.4.1 including
the XNNPACK delegate.}
\label{tab:external}
\centering
\small
\begin{tabular}{lrrrrr}
\toprule
 & \textbf{ours} & torch & torch & ET & ET \\
shape & \textbf{(V7)} & eager & compile & portable & XNNPACK \\
\midrule
\multicolumn{6}{l}{\emph{LayerNorm}} \\
1024$\times$2048 & \textbf{16.47} & 7.89 & 4.63 & 6.17 & 6.06 \\
4096$\times$8192 & \textbf{13.52} & 7.35 & 4.44 & 5.92 & 5.93 \\
\midrule
\multicolumn{6}{l}{\emph{RMSNorm}} \\
1024$\times$2048 & \textbf{38.45} & 3.44 & 6.65 & 4.27 & 6.64 \\
4096$\times$8192 & \textbf{24.63} & 2.77 & 5.90 & 3.26 & 4.19 \\
\bottomrule
\end{tabular}
\end{table}

\subsection{The result is about fusion, not kernel quality}

We are ahead of both baselines at every shape (Fig.~\ref{fig:baselines}):
LayerNorm by 1.5--2.1$\times$ and
RMSNorm by 4.2--5.8$\times$. The asymmetry is the interesting part, and it is
not about our kernels. Profiling what each module dispatches to on CPU shows
\code{torch.nn.LayerNorm} reaching a single fused kernel, while
\code{torch.nn.RMSNorm} decomposes into elementwise multiply, power, sum and
divide---each pass reading and writing the whole tensor. A fused RMSNorm exists
at the dispatcher level, but CPU eager still decomposes; compiling the module
fuses it and roughly doubles its throughput.

That produces a clean inversion. Within our implementation RMSNorm is
1.8--2.3$\times$ faster than LayerNorm, exactly as the traffic ratio says it
should be. Within PyTorch eager the ordering \emph{flips}, and its LayerNorm is
$2.7\times$ faster than its RMSNorm---purely because the cheaper operation is
the one without a fused CPU kernel. \emph{A fused implementation of the
expensive norm beats a decomposed implementation of the cheap one}, which is the
clearest argument in this report for why writing the kernel was worth doing.

\subsection{What this comparison does not establish}

Beating two general-purpose frameworks is not being state of the art, and we do
not claim it. No specialist vendor kernel was measured. We investigated Apple's
BNNS and obtained no number: on this OS its LayerNorm entry point is a
deprecated generic constructor we could create but not execute, and there is no
RMSNorm in BNNS at all---every ``RMS'' symbol in the headers is the optimizer.
Current Accelerate does expose both operations, but they post-date this OS
release. We therefore make \emph{no claim} about how our kernels compare with a
specialist vendor implementation, in either direction. An earlier draft did make
such a claim, citing a figure from our own error log that had never passed the
correctness-gated harness; it is withdrawn rather than restated.

Framework margins are also version statements. On the previous stack the same
harness would have made our RMSNorm margin read 8--13$\times$ instead of
4.2--5.8$\times$. The kernels did not change; the baseline did. Versions are
pinned and the run manifests are checked in.

% =========================================================================
\section{External Validation}

Every figure above comes from our own harness, which is a single-instrument
story, and a reader is entitled to ask whether the instrument is right. An
independent characterization of the same architecture agrees wherever the two
overlap: the streaming vector length matches exactly; their SSVE \code{FMLA}
throughput of 31\,GFLOP/s matches our independently derived 31.0\,GFLOPS, from a
different harness built for a different purpose; and their sharply reduced
bandwidth above $\approx$8\,MiB independently corroborates the footprint
correction of Sec.~\ref{sec:ceiling}, turning ``our probe produced these
numbers'' into ``our numbers reproduce an independently observed architectural
feature''. Their report that SME2 four-register loads reach $\approx$925\,GiB/s
against $\approx$376 for single-register loads independently motivates V7 as a
lever aimed at a documented feature rather than opportunistic ISA trickery.

Absolute bandwidth figures should \emph{not} be expected to match, and ours do
not: their measurements are read-only loads with a specific instruction
selection, ours a 1R+1W scale-add under the useful-bytes convention. Different
traffic definitions give different absolute numbers; the qualitative behaviour
is what agrees, and that is what is claimed.

% =========================================================================
\section{Threats to Validity}

\textbf{One machine, one OS.} Every number is from a single M4 on one macOS
release. Nothing here establishes behaviour on other SME implementations; the
kernels are vector-length-agnostic precisely because SVL is not guaranteed to be
512 bits.

\textbf{Best-case reporting, now bounded.} Headline figures are minima. Median
and p10--p90 are reported alongside and agree to $\approx$1\,\%, so the
best-case figures were typical---but that was an assumption until we measured
it.

\textbf{No specialist vendor baseline}, so no claim is made in either direction
(Sec.~\ref{sec:external}).

\textbf{The native-layout comparison answers one question, not two.} It does not
measure the cost of substituting our kernel into a framework's tensor layout.

\textbf{Framework margins are version statements}, and ExecuTorch figures
include per-invocation runtime dispatch---honest for an on-device comparison,
but not a pure kernel number.

\textbf{Single-threaded except where stated.} The frameworks parallelize by
default; our threaded results are reported separately against a chip ceiling.

\textbf{Group granularity distorts small shapes.} Below 64 rows the kernel does
a full group's work regardless, so small-$M$ figures reflect that rather than
bandwidth.

\textbf{One mechanism remains unidentified}: the base-offset sensitivity behind
the threading cliff (Sec.~\ref{sec:jit}). The mitigation is justified by a
control, not by an explanation.

% =========================================================================
\section{Related Work}

LayerNorm~\cite{ba2016layernorm} and RMSNorm~\cite{zhang2019rmsnorm} define the
operators; the latter reports comparable task quality at reduced cost, a
modelling claim our kernel-level measurements neither support nor contradict.
The roofline model~\cite{williams2009roofline} frames the analysis, and our
central methodological correction is essentially a statement that the roofline's
memory ceiling is itself a function of working-set size, which is standard in
principle and easy to violate in practice. Production CPU implementations of
LayerNorm commonly use Welford's recurrence to escape the naive formulation;
our measurements agree that Welford is far better than naive while being
comparable to, not dominated by, the centred two-pass---and slower on SIMD for a
reason specific to vectorization. Arm's architecture reference
manual~\cite{armarm} is the authority for SME/SSVE encodings, and we verified
every encoder against toolchain-assembled golden words rather than against
recollection.

% =========================================================================
\section{Conclusion}

We took both Transformer normalization primitives from a scalar reference to
JIT-generated SME kernels invoked through a compiler runtime, verifying each
step against a float64 oracle. The best kernels reach $\approx$33\,\% of a
footprint-matched ceiling on useful bytes ($\approx$50\,\% on moved bytes) at
cache-resident shapes and $\approx$42\,\% at DRAM, and beat two general-purpose
frameworks by 1.5--5.8$\times$.

The more useful outcomes are the negative and corrective ones. The ZA tile is
the wrong tool for a bandwidth-dominated horizontal reduction, and it took a
measurement---twice, since our first reasoned skip rested on an untested
factor---to say so. SME2 helps, but not for either reason we predicted. The
streaming transition is 9\,ns rather than the dominant per-call cost, and the
real small-tensor penalty is a granularity decision of our own making. And our
own roofline was mis-stated for five sprints in a way that inflated our results
by about 2$\times$.

The pattern behind every one of those errors is the same, and it is the finding
we would carry to the next project: each came from reasoning about a mechanism
instead of measuring it. The claims that survived contact were the ones written
down as a prediction \emph{before} the measurement, which is also the only
reason the errors were caught at all.

Scope note: SME was course material; SME2-specific multi-vector optimization was
not, and was reached here by detecting the hardware's actual feature set and
then testing a hypothesis against it.

% =========================================================================
\appendices

\section{Reproducibility}

\begin{verbatim}
cmake --preset release-submission
cmake --build build-submission -j
./build-submission/src/norm/main_norm
\end{verbatim}

The benchmark prints its own provenance header---git commit (marked
\code{-dirty} on a modified tree), build type, compiler, OS, CPU, detected
\code{FEAT\_SME}/\code{FEAT\_SME2}, streaming vector length, dispatch target,
thread count---so no table can be separated from the conditions that produced
it. The captured output behind every figure and table here is checked into the
repository, as are the pinned external-baseline environment and the run
manifests. Correctness runs with \code{ctest}; the norm suite is 137 cases and
780{,}264 assertions, and passes cleanly under AddressSanitizer and
UndefinedBehaviorSanitizer. Continuous integration builds every kernel and runs
the host-portable tests, but the runners have no SME: green CI cannot execute a
single SME kernel, and SME coverage is local to the target machine.

\section*{Acknowledgment}

We thank the Machine Learning Compilers Lab at Friedrich-Schiller-Universit\"at
Jena. Generative AI assistance was used for implementation support, for review
of our claims and measurements against the hardware, and for drafting text that
we subsequently rewrote; it is disclosed in full in the project repository.

\begin{thebibliography}{9}

\bibitem{ba2016layernorm}
J.~L. Ba, J.~R. Kiros, and G.~E. Hinton, ``Layer normalization,''
\emph{arXiv:1607.06450}, 2016.

\bibitem{zhang2019rmsnorm}
B.~Zhang and R.~Sennrich, ``Root mean square layer normalization,'' in
\emph{Advances in Neural Information Processing Systems (NeurIPS)}, 2019.

\bibitem{vaswani2017attention}
A.~Vaswani \emph{et al.}, ``Attention is all you need,'' in
\emph{Advances in Neural Information Processing Systems (NIPS)}, 2017.

\bibitem{touvron2023llama}
H.~Touvron \emph{et al.}, ``LLaMA: Open and efficient foundation language
models,'' \emph{arXiv:2302.13971}, 2023.

\bibitem{williams2009roofline}
S.~Williams, A.~Waterman, and D.~Patterson, ``Roofline: An insightful visual
performance model for multicore architectures,'' \emph{Communications of the
ACM}, vol.~52, no.~4, pp.~65--76, 2009.

\bibitem{armarm}
Arm Ltd., \emph{Arm Architecture Reference Manual for A-profile architecture},
SME and Streaming SVE chapters.

\bibitem{pytorch}
A.~Paszke \emph{et al.}, ``PyTorch: An imperative style, high-performance deep
learning library,'' in \emph{NeurIPS}, 2019.

\bibitem{executorch}
Meta AI, ``ExecuTorch: On-device inference runtime,''
\url{https://pytorch.org/executorch/}.

\bibitem{mccalpin1995stream}
J.~D. McCalpin, ``Memory bandwidth and machine balance in current high
performance computers,'' \emph{IEEE TCCA Newsletter}, pp.~19--25, 1995.

\end{thebibliography}

\end{document}
